\documentclass[../main.tex]{subfiles}
\begin{document}

\chapter{Edition of the Arabic Text}\label{chap:edition}

The following edition is based on \witness{P} and \witness{I}.
The variant at line \edlineref{ed:firaq} is discussed below.

\begin{Arabic}
\beginnumbering
\pstart
قال \edtext{جالينوس}{\Afootnote{جالينس \witness{I}}} إن
\edtext{الفرق}{\edlabel{ed:firaq}\Afootnote{الفرقة \witness{P}}\Bfootnote{Gal. \textit{De sectis} 1}}
في الطب ثلاث.
\pend
\pstart
\edtext{وأما أصحاب \edtext{التجربة}{\Afootnote{التجارب \witness{I}}} فإنهم يقولون}{\lemma{وأما \ldots{} يقولون}\Afootnote{om. \witness{P}}}
إن الطب يستخرج بالتجربة.
\pend
\endnumbering
\end{Arabic}

\section{An English edition}

\beginnumbering
\pstart
Galen says that the \edtext{sects}{\Afootnote{schools \witness{P}}} of
medicine are \edtext{three}{\Afootnote{two \witness{I}}\Bfootnote{Cf.
\cite[45]{Pormann2026}.}}, namely the Empiricist, the Dogmatist and the
\edtext{\textgreek{μεθοδικοί}}{\Afootnote{\textgreek{μεθοδικός} \witness{P}}}.
\pend
\endnumbering

A regular footnote after the edition.\footnote{An ordinary footnote.}

\end{document}
